\documentclass[11pt,a4paper]{article}
\usepackage[margin=2cm]{geometry}
\usepackage[T1]{fontenc}
\usepackage{enumitem,xcolor,parskip}
\usepackage{amssymb}
\setlist{nosep,leftmargin=1.2em}
\newcommand{\tbd}[1]{\textcolor{red}{[#1]}}
\pagestyle{empty}
\begin{document}
\begin{center}{\Large\bfseries Information and Consent Form}\\[2pt]
SeeSound pronunciation recording study --- DRAFT v0.2\end{center}

\textbf{Who we are.} A research team at Ho Chi Minh City University of Foreign Languages - Information Technology (HUFLIT) and National Formosa University. Contact: Prof.\ Hui-Kai Su, hksu@nfu.edu.tw. The data are managed by Kiet Pham (HUFLIT).

\textbf{Why we ask.} We are testing a web program (SeeSound) that shows deaf and hard-of-hearing learners how to say English consonants. Before it is used with children, we want to check, with teachers, how well the program's judgements of a spoken word agree with what teachers hear. This is a small pilot.

\textbf{What happens.} You (or your child) say eight English words (\textit{she, see, fee, pea, bee, me, knee, cheat}) twice each into a computer microphone. This takes about 4 minutes. The screen does not show a score. A teacher may later listen to the recordings and rate them.

\textbf{What we keep.} A participant code (not a name), the recordings \emph{only if you tick the box below}, and numbers measured from them. No camera or video is used. Data stay on a research computer managed by Kiet Pham; recordings are heard only by the raters and the research team, are not sent by e-mail or chat, and are deleted after 1 year. Reports show combined results only.

\textbf{Risks and benefits.} No known risks beyond a few minutes of speaking. There is no payment and no direct benefit; the study may help improve tools for learners.

\textbf{Your rights.} Taking part is voluntary. You or your child may stop at any time, and you may ask us to delete the recordings and data, without giving a reason and without any effect on your teaching or care. To withdraw, give the participant code and contact us at the address above.

\medskip\hrule\medskip
\textbf{Consent} (for a child, the parent or guardian completes this; the child is also asked whether they are happy to take part)

\begin{itemize}[label=$\square$]
\item I have read (or had read to me) this information and could ask questions.
\item I agree that I / my child takes part in the recording described above.
\item I \textbf{also} agree that the \textbf{audio recordings} are kept and listened to by the raters and research team as described. \textit{(If this is not ticked, the person can still take part, but teachers cannot rate the recordings.)}
\end{itemize}

\vspace{6pt}
\newcommand{\line}[2]{\noindent #1 \hrulefill\hspace{1em}#2 \rule{4cm}{0.4pt}\\[12pt]}
\noindent Participant code (filled in by the teacher): \rule{3.5cm}{0.4pt}\hfill Date: \rule{4cm}{0.4pt}\\[12pt]
\noindent Name of adult participant or parent/guardian: \hrulefill\\[12pt]
\noindent Signature: \hrulefill\\[12pt]
\noindent Teacher's initials (who explained the study): \rule{3cm}{0.4pt}\\[12pt]
\noindent Child's assent (``I am happy to do this''), if able: \rule{4cm}{0.4pt}

\vspace{4pt}\footnotesize Keep this form separate from the recordings: the form has a name, the data have only the code.
\end{document}
